\section{Introduction}
\label{sec:intro}

Every hospital transfer, every discharge, every clinical handoff carries a risk: a
medication gets dropped from the list, a dose gets duplicated, or a drug a patient
has been taking for years simply disappears from the record. Medication reconciliation
is the process of comparing a patient's complete medication history against every new
prescription, and it exists specifically to catch these errors. Medication discrepancies
occur in over 60\% of patients at care transitions, and a significant fraction of these
discrepancies lead to preventable adverse events \citep{asgari2025clinicalsafety}.

Large language models are beginning to be proposed and deployed for medication
reconciliation, often as part of broader FHIR-based EHR query pipelines. The appeal
is clear: LLMs can read unstructured clinical notes, understand medication names and
dosage semantics, and produce structured output, capabilities that rule-based systems
often struggle with. But the standard format for patient health records, FHIR R4,
is deeply nested JSON built around numeric ontology codes (RxNorm, SNOMED, LOINC)
rather than readable medication names. A single patient's complete FHIR bundle can
exceed 800,000 lines of JSON \citep{schmiedmayer2025llmonfhir}, far beyond any current
context window. Before a model can reason over a patient's medication history, that
history must be preprocessed, selected, and formatted into a prompt.

Prior work has recognised the preprocessing challenge. \citet{schmiedmayer2025llmonfhir}
and \citet{schmiedmayer2024llmonfhir} use function-calling to select relevant FHIR
resource types before passing data to a model. \citet{wu2026plannerauditor} apply
deterministic auditing pipelines on top of LLM-generated FHIR plans. What none of
these works have systematically studied is a more fundamental question: given a
task-relevant subset of FHIR data, does the \textit{format} in which that data is
presented to the model matter? Should it be raw JSON? A markdown table? A clinical
narrative? A chronological timeline?

This is not a hypothetical question. Every integration pipeline that sends FHIR data
to an LLM makes an implicit serialisation choice. That choice is almost always made
based on convenience or convention, not on empirical evidence about which format
produces better outcomes. We set out to generate that evidence.

\paragraph{Contributions.} This paper makes the following contributions:

\begin{itemize}
  \item A controlled benchmark of 200 synthetic patients $\times$ 4 serialisation
        strategies $\times$ 5 open-weight models = 4,000 total inference runs, measuring
        precision, recall, and F1 for medication name extraction from FHIR bundles.

  \item Evidence that serialisation strategy has a statistically significant,
        clinically meaningful effect on performance for models up to 8B parameters
        (Mistral-7B: $r = 0.617$, $p < 10^{-10}$), and that the optimal strategy is
        model-size-dependent: Clinical Narrative is best for $\leq$8B models; Raw JSON
        is best at 70B.

  \item Evidence that omission dominates hallucination in this task. Mean precision
        exceeds mean recall in all 20 model and strategy combinations, though
        hallucinations do occur, which changes how clinical safety auditing of these
        systems should be framed.

  \item Evidence of a hard capacity ceiling for smaller models: recall degrades sharply
        for patients with more than 7--10 concurrent active medications, affecting the
        highest-risk polypharmacy patients.

  \item A reproducible open pipeline using Synthea, Ollama, and no proprietary APIs,
        running on an AWS \texttt{g6e.xlarge} instance (NVIDIA L40S, 48 GB VRAM).
\end{itemize}

The rest of this paper is structured as follows. Section~\ref{sec:related} reviews
related work; Section~\ref{sec:methods} describes the dataset, strategies, models,
and evaluation; Section~\ref{sec:results} reports findings; Section~\ref{sec:discussion}
interprets them for clinical deployment; Section~\ref{sec:limitations} discusses
limitations.
